微分方程交到手上的时候，它只肯说一件事：此刻的变化率是多少。报告里要的却是整条轨迹。把前者变成后者，办法说穿了只有一个——站在当前点上读斜率，沿着它走一小段，把落点当作新的当前点，重复。整个单元都在追问：这样走一步，要拿什么去换。

要付出的分三处。第一处是精度：每一步的直行都略微偏离真曲线，走完全程之后这些偏离会累起来。第二处是稳定性：步长大到某个门槛，数值解会做出真解从未有过的行为——本该衰减的东西开始增长，而且这个门槛常常由一个你并不关心的快速分量定下。第三处出现在边值问题里：条件分散在区间两端，逐步推进的框架失去了出发点。

\begin{center}
\begin{tikzpicture}[x=1cm,y=1cm,font=\footnotesize]
  \foreach \i/\ta/\tb in {0/{沿斜率走一步}/{局部 \(O(h^2)\)，全局 \(O(h)\)},
                          1/{一步内多探几次}/{用函数求值换阶数},
                          2/{步长被稳定性绑架}/{显式的圆盘，隐式的半平面},
                          3/{两端都有约束}/{打靶还是全局离散}}{
    \draw[rounded corners=2pt,fill=black!5,draw=black!45]
      (\i*3.85,0) rectangle (\i*3.85+3.45,1.35);
    \node[align=center,text width=3.2cm] at (\i*3.85+1.725,0.68) {\ta\\[1pt]\tb};
  }
  \foreach \i in {0,1,2}{
    \draw[->,>=stealth,black!60] (\i*3.85+3.5,0.68) -- (\i*3.85+3.8,0.68);
  }
  \node[black!70] at (3.65,-0.45) {误差要结账};
  \node[black!70] at (7.5,-0.45) {快模态逼小步长};
  \node[black!70] at (11.35,-0.45) {无法单向推进};
\end{tikzpicture}

\emph{四篇的主线：每往右一格，都是前一格暴露出的困难逼出来的下一步。}
\end{center}

先问每种方法在赌什么。Euler 赌一步之内斜率变化不大；Runge--Kutta 不赌，它在一步内多问几个地方再把答案加权平均；隐式方法愿意每步解一个方程，换来对快速衰减分量的强阻尼。第三篇要说明的是，稳定性和精度是两把不同的尺子，量出来的步长上限可以差好几个数量级；被卡住的那一头往往不是你以为的那一头。第四篇则要求把``从左往右走''的习惯放下，改成把整条曲线的未知值一次联立。

这些判断在 AI 工程里有直接的对应物。学习率被 Hessian 最大特征值卡住、自适应优化器按方向调节步幅、扩散模型的采样步数与采样器阶数之争，都是同一套语言的另一种说法。

前一单元：\hyperref[ch:065]{``数值线性代数：结构、规模与保证的强弱决定算法''}。

\subsection{本章篇目}

以下篇目按概念依赖排列；若从具体问题进入，也可以先打开最接近当前困惑的一篇，再沿篇内链接回补。

\begin{itemize}
\tightlist
\item \hyperref[sec:08-Numerical-Analysis/06-ODE/01]{``Euler 方法：从局部斜率到全局误差''}；
\item \hyperref[sec:08-Numerical-Analysis/06-ODE/02]{``Runge--Kutta 与自适应步长''}；
\item \hyperref[sec:08-Numerical-Analysis/06-ODE/03]{``稳定区域与刚性方程''}；
\item \hyperref[sec:08-Numerical-Analysis/06-ODE/04]{``边值问题：从打靶到全局离散''}。
\end{itemize}

若时间只够读一篇，读第三篇：调不好的求解器，绝大多数卡在稳定性而不是精度上。
